\section{模型的建立与求解}
\subsection{问题一的模型建立与求解}
\subsubsection{问题一模型的建立}
\subsubsection*{step1 径向传热传质模型的建立}
药材初始半径$R_0=0.02$ m、长度$L=0.25$ m，固定域为$0\le r\le R_0$。设温度场为$T(r,t)$、干基含水率场为$C(r,t)$。附件1每60 s给出一组环境数据，覆盖0--14400 s；第一问使用其中0--1800 s的部分。对任一环境量$a$，在相邻观测之间令
\begin{equation}
a(t)=a_\ell+\frac{a_{\ell+1}-a_\ell}{t_{\ell+1}-t_\ell}(t-t_\ell),
\qquad t_\ell\le t\le t_{\ell+1}.
\label{eq:air}
\end{equation}

我们在该处采用分段线性插值。若用高阶多项式拟合这几十个数据点，容易引发著名的龙格现象，在数据点之间产生剧烈的虚假上下震荡，导致数值仿真失真；而分段线性插值单调平滑、严格保真，数值最稳定。环境温度$T_a(t)$与有效平衡含水率$C_e(t)$均按式\eqref{eq:air}处理，它们是已知边界输入，内部状态则需由控制方程求出。

\fig{fig1}{0.89}{圆柱药材径向传递示意：热量从环境传入，水分向外迁移}{fig:geometry}

取半径$r$至$r+\dd r$间的薄壳，其侧面积与体积主项为
\begin{equation}
A(r)=2\pi rL,\qquad
\dd V=\pi L[(r+\dd r)^2-r^2]\simeq2\pi rL\dd r .
\label{eq:shell}
\end{equation}

同样厚度的薄壳在不同半径处具有不同面积，故守恒时必须比较面积乘通量。先以干物质量为水分基准，有
\begin{equation}
C=\frac{m_w}{m_d},\qquad m_w=Cm_d,\qquad
m_d=\rho_d\dd V .
\label{eq:drybasis}
\end{equation}

在参考干密度固定时，薄壳内水分积累率即$\rho_d C_t\dd V$。温度方程则通过材料质量$\rho\dd V$与比热容$c_p$把温升转换为显热变化。两种储存量的基准不同，却可通过局部传递定律进入同一径向几何框架。

\subsubsection*{step2 从能量收支建立温度方程}
薄壳温度在短时间$\dd t$内变化约为$T_t\dd t$。第一问$\rho,c_p$固定，因此显热增量写为
\begin{equation}
\dd E\simeq \rho c_pT_t\dd V\dd t .
\end{equation}

规定径向热流密度$q$向外为正。薄壳内侧流入、外侧流出之差为$[A(r)q(r)-A(r+\dd r)q(r+\dd r)]\dd t$，故有
\begin{equation}
\rho c_pT_t\dd V=A(r)q(r)-A(r+\dd r)q(r+\dd r).
\end{equation}

两边除以体积并令薄壳厚度趋于零，代入傅里叶定律$q=-kT_r$，得到
\begin{equation}
\rho c_pT_t=-\frac1r\frac{\pd(rq)}{\pd r}
=\frac1r\frac{\pd}{\pd r}\left(rk\frac{\pd T}{\pd r}\right).
\label{eq:heat}
\end{equation}

式\eqref{eq:heat}左端是单位体积的显热积累率，右端是净导热输入率，二者量纲均为W/m$^3$。其中$1/r$来自圆柱侧面积随半径变化，并非附加热源。第一问采用题给常数，可展开为
\begin{equation}
T_t=\alpha\left(T_{rr}+\frac{T_r}{r}\right),\qquad
\alpha=\frac{0.36}{820\times2600}
\approx1.6886\times10^{-7}\ \mathrm{m^2/s}.
\label{eq:linearheat}
\end{equation}

温度的初值、中心对称条件和表面牛顿冷却条件分别为
\begin{align}
T(r,0)&=28,\qquad T_r(0,t)=0,\\
-kT_r(R_0,t)&=h[T_s(t)-T_a(t)],\qquad h=25\ \mathrm{W/(m^2K)}.
\label{eq:heatbc}
\end{align}

空气比表面热时，右端为负，说明热量向内传递。也可将边界写成$kT_r=h(T_a-T_s)$，两种表达等价。有限换热系数允许空气与表面之间存在温差，因此不能再附加$T_s=T_a$，否则会重复且一般不相容地约束表面。

\subsubsection*{step3 从水质量守恒建立非线性扩散模型}
采用有效菲克关系，将相对骨架的水分通量写为
\begin{equation}
j=-\rho_dD(C)C_r,\qquad
D(C)=7\times10^{-9}\exp(-0.89/C)\ \mathrm{m^2/s}.
\label{eq:fick}
\end{equation}

由于$C$是干基质量比，$D C_r$的单位为m/s；乘以$\rho_d$后才得到kg/(m$^2$s)的质量通量。对薄壳应用水质量守恒，并利用固定、均匀参考干密度，可得
\begin{align}
\rho_d C_t&=-\frac1r\pd_r(rj)
=\frac1r\pd_r(r\rho_dD(C)C_r),\\
C_t&=\frac1r\pd_r[rD(C)C_r].
\label{eq:water}
\end{align}

即使第一问热物性不变，水分方程仍是非线性的。展开式为
\begin{equation}
C_t=D(C)\left(C_{rr}+\frac{C_r}{r}\right)
+\frac{0.89}{C^2}D(C)(C_r)^2 .
\label{eq:nonlinearwater}
\end{equation}

最后一项来自对$D(C(r,t))$求空间导数，不能删除。程序不直接计算这些展开项，而通过共享界面通量保留式\eqref{eq:water}的守恒结构。

表面传质取$j_s=\rho_{d,s}h_m(C_s-C_e)$，与内部通量连续，约去相同的表面干密度后得到
\begin{align}
C(r,0)&=2.55,\qquad C_r(0,t)=0,\\
-D(C_s)C_r(R_0,t)&=h_m[C_s-C_e(t)],
\quad h_m=8\times10^{-7}\ \mathrm{m/s}.
\label{eq:waterbc}
\end{align}
本题没有给出材料解吸等温线，因此按假设4使用附件水分数值作为有效$C_e(t)$，以此与题给$h_m$共同闭合边界。

式\eqref{eq:heat}、\eqref{eq:water}及相应初边值条件构成第一问完整模型。第一问$D$仅依赖$C$，热物性均为常数，所以两个方程在物理上解耦；统一放入求解器只是数值实现的组织方式，并未人为添加温湿反馈。

初始状态下$D(2.55)\approx4.94\times10^{-9}$ m$^2$/数量级显著小于热扩散率，水分迁移远慢于温度调整。同时，$C$下降会进一步降低$D$，使表面失水区域的水分扩散能力持续减弱。二者共同导致表面附近形成含水率梯度陡峭的薄层，因而径向网格在表面附近需要较高分辨率，这是后续模型求解中采用表面加密网格的直接依据。

\subsubsection*{step4 表面加密网格与有限体积离散}
令$x=r/R_0$，把径向区间归一化到$[0,1]$。对$N$个区间定义
\begin{equation}
u_i=\frac{i}{N},\qquad
x_i=1-\frac{\sinh[\beta(1-u_i)]}{\sinh\beta},
\quad i=0,\ldots,N .
\label{eq:mesh}
\end{equation}

取$\beta=5$，节点靠近表面时逐渐加密。主计算$N=2560$，含2561个节点。令内部控制面$x_{i+1/2}=(x_i+x_{i+1})/2$，并令$x_{-1/2}=0$、$x_{N+1/2}=1$。相应控制体积的径向权重为
\begin{equation}
\omega_i=\int_{x_{i-1/2}}^{x_{i+1/2}}x\dd x
=\frac{x_{i+1/2}^2-x_{i-1/2}^2}{2},
\qquad \sum_{i=0}^{N}\omega_i=\frac12 .
\end{equation}

在界面平均状态$(C_i+C_{i+1})/2$、$(T_i+T_{i+1})/2$处计算物性。\(H_{i+1/2}\) 是热通量（对应傅里叶定律 \(q = -k \nabla T\)）。它表示从第 \(i\) 个微元流向第 \(i+1\) 个微元的热量速率，由导热系数 \(k\) 和相邻节点的温度梯度 \((T_{i+1} - T_i)\) 决定。\(W_{i+1/2}\) 是水分质量通量（对应菲克定律 \(J = -D \nabla C\)）。它表示从第 \(i\) 个微元流向第 \(i+1\) 个微元的水分质量速率，由扩散系数 \(D\) 和相邻节点的含水率梯度 \((C_{i+1} - C_i)\) 决定。
\begin{equation}
\begin{aligned}
H_{i+1/2}&=-\frac{x_{i+1/2}k_{i+1/2}}{R_0^2}
\frac{T_{i+1}-T_i}{x_{i+1}-x_i},\\
W_{i+1/2}&=-\frac{x_{i+1/2}D_{i+1/2}}{R_0^2}
\frac{C_{i+1}-C_i}{x_{i+1}-x_i}.
\end{aligned}
\label{eq:fluxes}
\end{equation}

积分控制方程并以节点状态代表该控制体，得到
\begin{equation}
\frac{\dd T_i}{\dd t}
=\frac{H_{i-1/2}-H_{i+1/2}}{\omega_i\rho_i c_{p,i}},
\qquad
\frac{\dd C_i}{\dd t}
=\frac{W_{i-1/2}-W_{i+1/2}}{\omega_i}.
\label{eq:ode}
\end{equation}

中心面通量取零；表面控制面直接使用交换关系：
\begin{equation}
H_{-1/2}=W_{-1/2}=0,\quad
H_{N+1/2}=\frac{h(T_N-T_a)}{R_0},\quad
W_{N+1/2}=\frac{h_m(C_N-C_e)}{R_0}.
\label{eq:facebc}
\end{equation}

表面通量中的$1/R_0$与内部通量中的$1/R_0^2$不同，前者由归一化梯度与对流边界联合消去后得到。中心控制体具有有限体积，使用零中心通量可直接更新，避免在$r=0$处计算$T_r/r$。

该离散继承有限体积法的局部收支结构\cite{patankar}：同一界面通量以相反符号进入相邻单元，总和时内部交换自动抵消。非均匀节点之间使用真实间距，不能以统一的$1/N$代替。界面物性取平均状态是一种离散近似，其可信度由后文网格与独立解检验确认。

\nextpage
\subsubsection{问题一的模型求解}
把节点状态交错组成$\bm y=(T_0,C_0,T_1,C_1,\ldots,T_N,C_N)^\mathsf T$，离散系统写为
\begin{equation}
\frac{\dd\bm y}{\dd t}=\bm F(t,\bm y).
\end{equation}

表面细网格使局部扩散时间尺度很短，同时烘干观测跨度较大，系统具有刚性。显式积分格式为满足稳定性将被迫采用极小步长，计算效率极低。为此，本文采用 SciPy 的变阶 BDF（Backward Differentiation Formula）方法\cite{scipy,shampine,bdfdocs}，以隐式方式处理扩散约束。其一阶形式为
\begin{equation}
\bm y_{n+1}=\bm y_n+\Delta t\,\bm F(t_{n+1},\bm y_{n+1}),
\end{equation}

实际实现自适应选择1--5阶并使用数值微分公式修正。右端依赖相邻节点，雅可比矩阵呈稀疏带状；程序以复步长法求导\cite{martins}：
\begin{equation}
J_{ij}\approx
\frac{\operatorname{Im}F_i(t,\bm y+\mathrm{i}\varepsilon\bm e_j)}
{\varepsilon},\qquad \varepsilon=10^{-25}.
\label{eq:complex}
\end{equation}

对相隔三个节点的同类状态可同时扰动，采用六组着色减少求导次数。复数仅用于右端的导数评估，物理解仍保持实数；该方法要求受扰动运算局部解析，最大值事件不参与此求导过程。

\begin{figure}[H]\centering
\begin{tikzpicture}[node distance=7mm]
\node[box,text width=106mm] (a) {读环境并建网格：$T_i=28,\ C_i=2.55$};
\node[box,text width=106mm,below=of a] (b) {取当前环境；计算节点热容量、界面物性和通量};
\node[box,text width=106mm,below=of b] (c) {BDF隐式步：稀疏雅可比、非线性迭代、误差控制};
\node[box,text width=106mm,below=of c] (d) {达到1800 s；时间连续输出与空间PCHIP插值};
\draw[arr](a)--(b);\draw[arr](b)--(c);\draw[arr](c)--(d);
\draw[arr](c.east)--++(9mm,0)|-(b.east);
\end{tikzpicture}
\caption{问题一的计算流程；回路表示自适应时间推进}\label{fig:algorithm}
\end{figure}

第一问相对容差为$2\times10^{-11}$、最大步长2 s，绝对容差对温度和含水率分别按$0.1\,\mathrm{rtol}$、$0.01\,\mathrm{rtol}$设置。时间采用求解器连续输出，空间采用PCHIP保形三次插值\cite{fritsch,pchipdocs}，提取题目要求的21个半径位置。每秒输出不要求内部步长恰好为1 s，计算网格也不等同于输出位置网格。

\begin{table}[H]\centering
\caption{问题1：温度分布（$^\circ$C）}\label{tab:q1T}
\begin{tabular}{cccccc}\toprule
时间/s & 0 cm & 0.5 cm & 1 cm & 1.5 cm & 2 cm\\ \midrule
100 & 28.0001 & 28.0003 & 28.0040 & 28.0327 & 28.1801\\
300 & 28.0408 & 28.0635 & 28.1514 & 28.3681 & 28.8489\\
600 & 28.4533 & 28.5360 & 28.8039 & 29.3159 & 30.1652\\
900 & 29.3243 & 29.4583 & 29.8755 & 30.6161 & 31.7304\\
1200 & 30.5427 & 30.7098 & 31.2223 & 32.1126 & 33.4276\\
1500 & 31.9957 & 32.1867 & 32.7660 & 33.7463 & 35.1203\\
1800 & 33.5753 & 33.7720 & 34.3642 & 35.3621 & 36.7856\\
\bottomrule\end{tabular}\end{table}

\fig{fig3}{0.95}{问题一各代表时刻的径向温度分布}{fig:q1temp}

表\ref{tab:q1T}和图\ref{fig:q1temp}显示，在所考察时刻表面温度高于中心，径向温度由内向外平滑增加。1800 s时中心为33.5753 $^\circ$C，表面为36.7856 $^\circ$C，二者相差3.2102 $^\circ$C。同期空气温度为41.513 $^\circ$C，表面尚低于环境，说明外部换热阻力与内部热容量共同造成温度滞后。

在100 s时中心温度输出为28.0001 $^\circ$C，几乎保持初值，而表面已升至28.1801 $^\circ$C。这里的“四位小数几乎不变”是短时响应幅度小和显示精度共同作用的表现。

长度$L$在局部方程中约去，但在计算总热量时仍会进入体积和表面积。当前结果描述径向温度分布。第一问输出文件对应1--1800 s、0--2 cm的温度网格，其展示精度为四位小数。

\begin{table}[H]\centering
\caption{问题1：干基含水率分布（kg/kg）}\label{tab:q1C}
\begin{tabular}{cccccc}\toprule
时间/s & 0 cm & 0.5 cm & 1 cm & 1.5 cm & 2 cm\\ \midrule
100 & 2.5500 & 2.5500 & 2.5500 & 2.5500 & 2.2470\\
300 & 2.5500 & 2.5500 & 2.5500 & 2.5492 & 2.0508\\
600 & 2.5500 & 2.5500 & 2.5500 & 2.5353 & 1.8770\\
900 & 2.5500 & 2.5500 & 2.5497 & 2.5045 & 1.7547\\
1200 & 2.5500 & 2.5500 & 2.5482 & 2.4646 & 1.6586\\
1500 & 2.5500 & 2.5499 & 2.5445 & 2.4206 & 1.5787\\
1800 & 2.5500 & 2.5497 & 2.5383 & 2.3755 & 1.5102\\
\bottomrule\end{tabular}\end{table}

\fig{fig4}{0.95}{问题一各代表时刻的径向干基含水率分布}{fig:q1water}

水分变化集中在外层。1800 s时表面含水率降至1.5102 kg/kg，中心仍为2.5499924 kg/kg，四舍五入显示2.5500；距中心1.5 cm处已降至2.3755 kg/kg，1 cm处仅下降到2.5383 kg/kg。该分布与内层到表面的有限迁移速度相符。

表面虽与低水分环境交换，但有限$h_m$使表面值不会立即等于$C_e$。随着外层失水，局部$D(C)$下降，形成较陡梯度，进一步限制内部水分外迁。因此，表面下降明显和中心暂时变化很小并不矛盾，也不说明水分方程失去守恒。

按圆柱体积权重计算，1800 s时平均含水率约为2.2935572 kg/kg。由于外层对应较大的体积权重，不能简单对五个半径处数值取算术平均。第一问结果保存为\texttt{result1.xlsx}的“温度”和“水分浓度”两张表；其中每秒21个位置，共1800个时间行。
